%%%%%%%%%%%%%%%%%%%%%%%%%%%%%%%%%%%%%%%%%%%%%%%%%%%%%%%%%%%%%%%%%%%%%%%%%%%%%%%
% Chapter 6: Conclusion
%%%%%%%%%%%%%%%%%%%%%%%%%%%%%%%%%%%%%%%%%%%%%%%%%%%%%%%%%%%%%%%%%%%%%%%%%%%%%%%

\chapter{Conclusion}
\label{ch:conclusion}

This chapter summarizes the key findings of our research on the Temporal DGNN model for video anomaly detection using the NiAD\_large\_graphs dataset and discusses directions for future work.

\section{Summary of Contributions}

This research has made several contributions to the field of video anomaly detection using graph-based representations:

\subsection{Node-Feature-Based Data Pipeline}

We developed a streamlined data processing pipeline that:
\begin{enumerate}
    \item Converts 30-frame surveillance videos into sequences of PyTorch Geometric graphs
    \item Employs YOLOv8 for robust vehicle detection (car, motorcycle, bus, truck)
    \item Implements centroid-based tracking for consistent vehicle identification
    \item Produces dynamic graphs with variable node counts per frame
    \item Eliminates the need for padding, masks, and sentinel values
    \item Stores graphs efficiently as \texttt{.pt} files (30 per video) with CSV manifests
\end{enumerate}

\subsection{Temporal DGNN Architecture}

We implemented an effective Temporal DGNN model that:
\begin{itemize}
    \item Processes temporal sequences of per-frame graphs
    \item Uses two GCN layers for spatial feature extraction
    \item Employs node ID sorting for temporal consistency
    \item Applies mean pooling for frame-level embeddings
    \item Uses GRU for temporal sequence modeling
    \item Achieves \textbf{89.32\% test accuracy} with \textbf{AUC-ROC of 0.81} on the NiAD\_large\_graphs dataset
\end{itemize}

\subsection{Experimental Analysis}

Our comprehensive experimental analysis revealed:
\begin{enumerate}
    \item The effectiveness of dynamic graph representations
    \item The benefits of data augmentation for anomalous class balancing
    \item The impact of Focal Loss on handling class imbalance
    \item The importance of node ID alignment for temporal consistency
    \item The challenges of anomaly detection with imbalanced test sets
\end{enumerate}

\section{Key Findings}

\subsection{Architecture Insights}

\begin{enumerate}
    \item \textbf{Simplicity is effective:} The Temporal DGNNSimple architecture with just 2 GCN layers and a single GRU layer achieves strong performance.
    
    \item \textbf{Temporal modeling is crucial:} The GRU captures sequential dependencies that distinguish normal traffic from anomalous events.
    
    \item \textbf{Node ID alignment matters:} Sorting nodes by tracking IDs ensures consistent temporal feature alignment across frames.
\end{enumerate}

\subsection{Data Representation Insights}

\begin{enumerate}
    \item \textbf{Dynamic graphs are cleaner:} Variable-size per-frame graphs eliminate padding complexity and sentinel value handling.
    
    \item \textbf{Edge weights provide context:} Euclidean distances between vehicles encode spatial relationships effectively.
    
    \item \textbf{Augmentation helps balance:} 4$\times$ augmentation of anomalous samples improves training balance.
\end{enumerate}

\subsection{Practical Insights}

\begin{enumerate}
    \item \textbf{Class imbalance is critical:} Despite Focal Loss, the 8.4:1 test set imbalance limits anomaly detection recall.
    
    \item \textbf{High normal accuracy:} The model excels at identifying normal traffic patterns (F1 = 0.93).
    
    \item \textbf{Low anomaly recall:} Detecting anomalies remains challenging (recall = 13.64\%) due to class imbalance.
\end{enumerate}

\section{Limitations}

\subsection{Dataset Limitations}

\begin{enumerate}
    \item \textbf{Class imbalance:} The test set ratio of 8.4:1 (Normal:Anomalous) limits fair evaluation.
    
    \item \textbf{Limited anomaly samples:} Only 22 anomalous test samples provide limited statistical power.
    
    \item \textbf{Single domain:} All videos are from traffic surveillance night-time scenarios.
\end{enumerate}

\subsection{Model Limitations}

\begin{enumerate}
    \item \textbf{Low anomaly recall:} Only 13.64\% of anomalies are detected, insufficient for safety-critical applications.
    
    \item \textbf{Detection dependency:} YOLO detection quality directly affects graph quality.
    
    \item \textbf{Tracking errors:} Tracking failures can disrupt temporal consistency.
\end{enumerate}

\section{Future Work}

Future research directions for GNN-based video anomaly detection include:

\begin{enumerate}
    \item \textbf{Class imbalance mitigation:} Developing better sampling strategies, threshold optimization, and cost-sensitive learning approaches to improve anomaly detection recall.
    
    \item \textbf{Advanced temporal modeling:} Exploring bidirectional recurrent networks, Transformer-based temporal attention, and specialized temporal graph networks.
    
    \item \textbf{Enhanced feature engineering:} Adding velocity, relative position, and detection confidence as additional node features.
    
    \item \textbf{Anomaly localization:} Extending binary classification to temporal and spatial localization of anomalous events.
    
    \item \textbf{Real-time deployment:} Optimizing inference speed for practical traffic monitoring applications.
    
    \item \textbf{Self-supervised approaches:} Learning normal traffic patterns without labels for unsupervised anomaly detection.
\end{enumerate}

\section{Broader Impact}

\subsection{Positive Impacts}

\begin{enumerate}
    \item \textbf{Road safety:} Automated accident detection enables faster emergency response
    \item \textbf{Traffic management:} Real-time anomaly detection supports intelligent transportation
    \item \textbf{Research baseline:} Our pipeline and model serve as baselines for future research
\end{enumerate}

\subsection{Considerations}

\begin{enumerate}
    \item \textbf{Privacy:} Video surveillance raises privacy concerns
    \item \textbf{Reliability:} Safety-critical systems require high reliability
    \item \textbf{Bias:} Models may exhibit bias based on training data
\end{enumerate}

\section{Final Remarks}

This research demonstrates that Graph Neural Networks provide an effective framework for video anomaly detection by representing vehicle interactions as dynamic graphs. The Temporal DGNN model achieves \textbf{89.32\% accuracy} with \textbf{AUC-ROC of 0.81} on the NiAD\_large\_graphs dataset, with strong performance on normal traffic detection (97.3\% specificity).

The key contributions of our work are:
\begin{itemize}
    \item A clean, node-feature-based pipeline producing PyG \texttt{.pt} files
    \item Dynamic graph representation without padding or masks
    \item Effective temporal modeling through GCN + GRU architecture
    \item Comprehensive analysis of class imbalance effects
\end{itemize}

While the results are promising, significant challenges remain in achieving the high anomaly recall rates required for safety-critical applications. Addressing class imbalance and improving anomaly detection recall are critical areas for future research.

